\documentclass[11pt,a4paper]{article}

\usepackage[margin=1in]{geometry}
\usepackage{booktabs}
\usepackage{longtable}
\usepackage{array}
\usepackage{xcolor}
\usepackage{listings}
\usepackage{amsmath}
\usepackage{amssymb}
\usepackage{enumitem}
\usepackage{natbib}
\usepackage{hyperref}

\hypersetup{
  colorlinks=true,
  linkcolor=blue!50!black,
  citecolor=blue!50!black,
  urlcolor=blue!50!black,
  bookmarksopen=true
}

\lstdefinestyle{pystyle}{
  language=Python,
  basicstyle=\ttfamily\small,
  keywordstyle=\color{blue!70!black}\bfseries,
  commentstyle=\color{gray!70!black}\itshape,
  stringstyle=\color{green!35!black},
  showstringspaces=false,
  breaklines=true,
  breakatwhitespace=false,
  columns=flexible,
  frame=single,
  rulecolor=\color{black!30},
  numbers=none,
  tabsize=4,
  xleftmargin=0.5em,
  xrightmargin=0.5em
}
\lstset{style=pystyle}

\newcommand{\pkg}{\texttt{pyxtabond2}}
\newcommand{\note}[1]{\par\medskip\noindent\fbox{\begin{minipage}{0.96\linewidth}\small\textbf{Note.} #1\end{minipage}}\medskip}
\newcommand{\warn}[1]{\par\medskip\noindent\fbox{\begin{minipage}{0.96\linewidth}\small\textbf{Caution.} #1\end{minipage}}\medskip}
\newcommand{\bigO}{O}

\title{\textbf{PyXtabond2 Methodological Note}\\[0.4em]\Large GMM Theory, the IFE-GMM Extension, and Validation Against Stata's \texttt{xtabond2}}
\author{Frederick Kahindo}
\date{Version 0.0.4\\ \today}

\begin{document}

\maketitle

\begin{abstract}
\noindent This note provides the full econometric derivation underlying
\pkg{}, together with a complete, primary-source-grounded account of its
validation. Section~\ref{sec:gmm} derives the standard dynamic panel GMM
estimator (Difference and System GMM), its one- and two-step weighting
matrices, the \citet{Windmeijer2005} finite-sample variance correction, and
the small-sample degrees-of-freedom corrections, each tied explicitly to the
corresponding construction in Stata's \texttt{xtabond2.mata}
\citep{Roodman2009a}. Section~\ref{sec:diag} derives the Sargan/Hansen,
Difference-in-Sargan/Hansen, and Arellano-Bond AR($\ell$) test statistics.
Section~\ref{sec:ife} derives \emph{IFE-GMM}, an original extension not
present in \texttt{xtabond2}: an iterative, PCA-based estimator for dynamic
panels contaminated by interactive fixed effects \citep{Bai2009}, combining
automatic factor-count selection \citep{BaiNg2002, AhnHorenstein2013}, a
profile-GMM defactoring loop \citep{HongSuJiang2022}, and a split-panel
jackknife bias correction \citep{DhaeneJochmans2015}. Section~\ref{sec:val} reports, in full, the findings of two validation
passes conducted against \texttt{xtabond2.mata} and a real Stata run,
alongside every confirmed-correct result and every documented, unresolved
limitation.
\end{abstract}

\newpage
\tableofcontents

\newpage
\section{Introduction}

\subsection{Scope}

The companion \emph{User Manual} documents every \pkg{} option practically:
what it does, its Stata equivalent, and how to use it. This note instead
derives \emph{why} each computation is what it is, at the level of detail
needed to audit the implementation against its econometric sources. It is
organized in three parts: (i) the standard dynamic panel GMM machinery
(Sections~\ref{sec:gmm}--\ref{sec:diag}), present in Stata's \texttt{xtabond2}
and reproduced by \pkg{}; (ii) IFE-GMM (Section~\ref{sec:ife}), the package's
original extension, absent from \texttt{xtabond2}; and (iii) a full
validation account (Sections~\ref{sec:val}--\ref{sec:limitations}) --
confirmed-correct results and unresolved limitations documented rather than
silently assumed away.

\subsection{Notation}

$i=1,\dots,N$ indexes cross-sectional units, $t=1,\dots,T$ time periods.
$y_{it}$ is the dependent variable, $x_{it}$ a $p\times1$ vector of
regressors (possibly including precomputed lags of $y$), $\beta$ the
$p\times1$ coefficient vector, $\eta_i$ an unobserved time-invariant effect,
$\varepsilon_{it}$ an idiosyncratic shock. $\Delta$ denotes the first-difference
operator, $Z_i$ the instrument matrix for unit $i$ (stacked across units as
$Z$), $W$ a GMM weighting matrix. $H$ denotes a GLS-type covariance-structure
matrix (Stata's \texttt{h()}). For IFE-GMM, $\lambda_i$ ($r\times1$) is a
factor-loading vector, $F_t$ ($r\times1$) a common-factor vector, with
$\Lambda=(\lambda_1,\dots,\lambda_N)'$ ($N\times r$) and
$F=(F_1,\dots,F_T)'$ ($T\times r$).

\section{The Standard Dynamic Panel GMM Estimator}
\label{sec:gmm}

\subsection{The problem: fixed effects in a dynamic panel}

Consider the workhorse dynamic panel model
\begin{equation}
y_{it} = \rho\, y_{i,t-1} + x_{it}'\beta + \eta_i + \varepsilon_{it}. \label{eq:model}
\end{equation}
Because $y_{i,t-1}$ is a function of $\eta_i$, pooled OLS on
\eqref{eq:model} is inconsistent, and the within/LSDV (fixed-effects)
estimator that removes $\eta_i$ by demeaning is itself inconsistent for
fixed $T$ as $N\to\infty$: demeaning induces a correlation of order
$\bigO(1/T)$ between the demeaned lag and the demeaned error that does not
vanish as $N$ grows \citep{Nickell1981}. Two transformations remove $\eta_i$
\emph{without} this problem, provided instruments are used for the
resulting correlation between the transformed regressor and the transformed
error:

\begin{description}[nosep]
  \item[First differences] $\Delta y_{it} = \rho\,\Delta y_{i,t-1} +
    \Delta x_{it}'\beta + \Delta\varepsilon_{it}$, $t=2,\dots,T$.
  \item[Forward orthogonal deviations (FOD)] \citep{ArellanoBover1995}
    \[
      y_{it}^{\mathrm{FOD}} = c_{it}\left(y_{it} - \frac{1}{T_{it}^{+}}\sum_{s>t} y_{is}\right),
      \qquad c_{it} = \sqrt{\frac{T_{it}^{+}}{T_{it}^{+}+1}},
    \]
    where $T_{it}^{+}$ is the number of available future observations for
    unit $i$ at $t$. Unlike differencing, FOD only uses \emph{future}
    observations, so a single missing period does not additionally destroy
    the transformed values of neighbouring periods -- the reason it is
    preferred for unbalanced panels with gaps. This is exactly
    \texttt{PanelData.get\_fod}'s construction
    (\texttt{data\_utils.py::\_fod\_group\_vectorized}).
\end{description}

Both transformations still leave $\Delta y_{i,t-1}$ (or its FOD analogue)
correlated with the transformed error, since $y_{i,t-1}$ is correlated with
$\varepsilon_{i,t-1}$, which enters $\Delta\varepsilon_{it}$. This motivates
instrumenting rather than direct estimation.

\subsection{Moment conditions and the instrument matrix}
\label{sec:moments}

Under sequential exogeneity/predeterminedness, $E[y_{is}\varepsilon_{it}]=0$
for $s<t$ (more generally, for any GMM-style variable in \texttt{gmm\_vars}),
so that
\[
  E\!\left[y_{i,t-2}\,\Delta\varepsilon_{it}\right] = 0,\quad
  E\!\left[y_{i,t-3}\,\Delta\varepsilon_{it}\right] = 0,\ \dots
\]
for $t=2,\dots,T$: each time period's differenced error is orthogonal to
\emph{every} level of the instrumenting variable lagged two or more periods.
This yields the classic Arellano-Bond "staircase" instrument matrix, one row
block per unit \citep[matching][their notation, adapted]{Windmeijer2005}:
\[
  Z_i = \begin{bmatrix}
    y_{i1} & 0 & 0 & 0 & 0 & 0\\
    0 & y_{i1} & y_{i2} & 0 & 0 & 0\\
    & & \ddots & & &\\
    0 & 0 & 0 & y_{i1} & \dots & y_{i,T-1}
  \end{bmatrix},
\]
with $T(T-1)/2$ instrument columns for a single GMM-style variable when every
lag from 1 is admissible -- the maximal case, appropriate for a
\emph{predetermined} variable (uncorrelated with future, but not
necessarily contemporaneous, errors). \pkg{}'s own default
(\texttt{lag=(2, None)}) instead starts at lag 2, the safer assumption for a
genuinely \emph{endogenous} variable such as the lagged dependent variable
itself, and yields $(T-2)(T-1)/2$ columns; the construction below is
identical either way, only the lag window differs. This is exactly what
\texttt{SystemGMMBuilder.\_fill\_block\_a\_standard}
(\texttt{gmm\_builder.py}) constructs: one column per (time period, lag)
pair, subject to the chosen lag window. Two generalizations \pkg{}
implements on top of this baseline (both mirroring \texttt{xtabond2}):

\begin{itemize}[nosep]
  \item \textbf{Lag windows} (\texttt{lag\_limits\_diff}/\texttt{GMMStyle.lag}):
    restricting to lags $\ell\in[l_{\min},l_{\max}]$ instead of "2 to
    maximum", reducing instrument count and often improving the Sargan/Hansen
    test's power \citep{Roodman2009b}.
  \item \textbf{Collapsing} (\texttt{collapse}/\texttt{GMMStyle.collapse}):
    summing across time periods within each lag distance, so a variable
    contributes $l_{\max}-l_{\min}+1$ columns instead of $\bigO(T^2)$.
    This is the single most effective lever against \emph{instrument
    proliferation}: with too many instruments relative to $N$, GMM can
    overfit the endogenous variable in finite samples and the Sargan/Hansen
    overidentification test loses power to detect genuine misspecification
    \citep{Roodman2009b} -- a suspiciously perfect Hansen $p$-value is a
    symptom of this, not reassurance.
\end{itemize}

\subsection{System GMM: adding the levels equation}
\label{sec:system}

Difference GMM performs poorly when $y_{it}$ is highly persistent: lagged
levels are then weak instruments for the differenced equation
\citep{BlundellBond1998}. Blundell-Bond/Arellano-Bover System GMM adds the
\emph{levels} equation
\[
  y_{it} = \rho\, y_{i,t-1} + x_{it}'\beta + \eta_i + \varepsilon_{it}
\]
instrumented by \emph{lagged differences} of the GMM-style variable, valid
under the additional "mean stationarity" assumption
$E[\Delta y_{i2}\,\eta_i]=0$ (initial deviations from the long-run mean are
uncorrelated with the fixed effect). \texttt{SystemGMMBuilder.build\_system\_instruments}
builds, per GMM-style variable and per unit, a $2T\times n_{\text{cols}}$
block:
\begin{itemize}[nosep]
  \item Rows $1,\dots,T$ (\emph{differenced}-equation instruments): entry
    $(t,\cdot)=y_{i,t-\ell}$ for each admissible lag $\ell$ -- Section~\ref{sec:moments}'s block.
  \item Rows $T+1,\dots,2T$ (\emph{levels}-equation instruments): entry
    $(T+t,\cdot)=\Delta y_{i,t-d}$, $d=l_{\min}-1$. In the standard
    (non-collapsed) style this is $T-1$ \emph{time-specific} columns, each
    isolating $\Delta y_{i,t-d}$ at exactly period $t$ and zero elsewhere
    (\texttt{\_fill\_block\_b\_standard}); collapsing sums these into a
    single shared column (\texttt{\_fill\_block\_b\_collapse}), the levels-equation
    analogue of collapsing the differenced block.
\end{itemize}
The two blocks are stacked, giving $Z_i$ a genuine $2T\times n_{\text{cols}}$
shape even though a single unit contributes at most $2T-2$ real (non-missing)
rows -- unused rows/instruments are naturally zeroed and, for instrument
columns that are identically zero across the whole sample, dropped
(\texttt{GMMEngine.\_\_init\_\_}'s \texttt{non\_zero\_cols} filter). Standard
IV instruments (\texttt{iv\_vars}/\texttt{IVStyle}, for strictly exogenous
variables) are appended analogously via \texttt{build\_iv\_instruments}: a
variable is used as its own instrument in the differenced and/or levels
equation, and -- following \texttt{xtabond2}'s convention exactly --
appears as a \emph{single} instrument column combining both equations
whenever it instruments both, rather than as two separate columns.

\subsection{One-step and two-step GMM}
\label{sec:onetwostep}

Stacking $Y$, $X$, $Z$ across all units and (for System GMM) both equations,
the GMM estimator for a weighting matrix $W$ is
\begin{equation}
  \hat\beta(W) = \left(X'ZWZ'X\right)^{-1}X'ZWZ'Y, \label{eq:gmmest}
\end{equation}
the sample analogue of minimizing the quadratic form
$\left(\frac1N\sum_i Z_i'\hat\varepsilon_i\right)'W\left(\frac1N\sum_i
Z_i'\hat\varepsilon_i\right)$ over $\beta$ \citep{Hansen1982}. \textbf{One-step}
GMM ($\hat\beta_1$) uses a weight matrix $W_1$ that does \emph{not} depend on
estimated parameters -- $W_1 = \left(\sum_g Z_g'HZ_g\right)^{-1}$ for a fixed
$H$ (Section~\ref{sec:hmatrix}); \texttt{GMMEngine.\_compute\_W1} builds
exactly this sum, group by group. \textbf{Two-step} GMM ($\hat\beta_2$) uses
the efficient weight matrix built from the one-step residuals,
\begin{equation}
  W_2 = \left(\sum_g Z_g'\hat\varepsilon_{1,g}\hat\varepsilon_{1,g}'Z_g\right)^{-1}
      = S_2^{-1}, \label{eq:w2}
\end{equation}
computed via the identity $Z_g'ee'Z_g=(Z_g'e)(Z_g'e)'$ for efficiency
(\texttt{\_accumulate\_ZeZe}). $W_2$ is consistent for the inverse of the
true moment covariance under heteroskedasticity/serial correlation of
unrestricted form -- this is what makes two-step GMM asymptotically
efficient, and also what makes its \emph{naive} standard errors unreliable
in finite samples (Section~\ref{sec:windmeijer}).

\subsection{The \texttt{h()} weighting-matrix family}
\label{sec:hmatrix}

$W_1$'s $H$ matrix models $\mathrm{Cov}(\Delta\varepsilon_{it},\Delta\varepsilon_{is})$
under the auxiliary assumption that $\varepsilon_{it}$ is i.i.d.\ (a
reasonable working assumption for constructing an efficient \emph{one-step}
weight, later made robust to violations at the two-step/robust stage).
Since $\Delta\varepsilon_{it}=\varepsilon_{it}-\varepsilon_{i,t-1}$,
\[
  \mathrm{Cov}(\Delta\varepsilon_{it},\Delta\varepsilon_{is}) =
  \begin{cases} 2\sigma^2 & s=t\\ -\sigma^2 & |s-t|=1\\ 0 & |s-t|\ge2,\end{cases}
\]
giving a tridiagonal $H$ with 2 on the diagonal and $-1$ on the first
off-diagonals. This is exactly \texttt{\_build\_H\_diff\_block}'s
construction -- matching the matrix Windmeijer's own paper describes for the
differenced-equation weight ("2's
on the main diagonal, $-1$'s on the first off-diagonals"). Stata's
\texttt{h()} option (\texttt{h}, default 3) selects among three variants,
generalized in \pkg{} by \texttt{\_build\_H(h, T\_span, orthogonal)}:

\begin{description}[nosep]
  \item[\texttt{h=1}] $H=I$: assumes no serial-correlation structure at all
    in the transformed error, i.e.\ does not model the covariance induced by
    differencing/FOD itself. Rarely appropriate, but available for
    comparability with software that defaults to it.
  \item[\texttt{h=2}] The tridiagonal $H$ above for the differenced block;
    for System GMM, the cross-covariance between the differenced and levels
    blocks is assumed zero (block-diagonal $H$).
  \item[\texttt{h=3}] (Stata's default) As \texttt{h=2}, but additionally
    models the \emph{true} cross-covariance between the differenced and
    levels equations. Since
    $\mathrm{Cov}(\Delta\varepsilon_{it},\varepsilon_{is})=\sigma^2$ if
    $s=t$ and $-\sigma^2$ if $s=t-1$ (zero otherwise), the cross-block has
    $+1$ on its main diagonal and $-1$ on one off-diagonal -- exactly
    \texttt{\_build\_H\_full\_system}'s
    \texttt{H\_full[i,T\_span+i]=1.0} / \texttt{H\_full[i2+1,T\_span+i2]=-1.0}
    assignments.
\end{description}
\texttt{h=2} and \texttt{h=3} coincide for pure Difference GMM (no levels
block exists to model a cross-covariance with). The FOD case uses an
analogous but distinct construction (\texttt{\_xform\_matrix}), reflecting
FOD's own weights rather than the simple differencing weights, and not
reproduced in full here.

\subsection{Standard errors and the Windmeijer (2005) correction}
\label{sec:windmeijer}

\subsubsection{One-step}

Non-robust one-step variance is $V_1=(X'ZW_1Z'X)^{-1}$, scaled by an
estimated error variance $\hat\sigma^2$ computed from the \emph{same}
equation $H$ models (Section~\ref{sec:sig2} details exactly which residuals
enter $\hat\sigma^2$ and why this depends on \texttt{h}). Robust one-step
variance uses the sandwich
\[
  V_{1,\mathrm{rob}} = V_1\left(X'ZW_1\right)S_2\left(W_1Z'X\right)V_1,
\]
with $S_2$ as in \eqref{eq:w2} -- \texttt{estimator.py}'s one-step-robust
branch. Additionally, when \texttt{robust=True} with one-step estimation,
\pkg{} replicates \texttt{xtabond2}'s "Roodman's trick": a silent two-step
refit computed \emph{only} to obtain a valid Hansen statistic (the
non-robust Sargan statistic is otherwise the only overidentification test
available at one-step), without altering the reported one-step
$\hat\beta$/se.

\subsubsection{Two-step: why naive standard errors are wrong}

$W_2$ in \eqref{eq:w2} is itself a function of $\hat\beta_1$ (through
$\hat\varepsilon_1$). The naive two-step variance
$V_2=(X'ZW_2Z'X)^{-1}$ treats $W_2$ as if it were the \emph{true,
non-stochastic} $W$, ignoring the extra sampling variability it inherits
from $\hat\beta_1$. \citet{Windmeijer2005} shows this omission is the
dominant source of the well-documented downward bias of two-step GMM
standard errors in finite samples. Writing the moment vector as
$\bar g(\beta)=\frac1N\sum_i g_i(\beta)$ and the efficient weight as
$W_N(\hat\beta_1)=\frac1N\sum_i g_i(\hat\beta_1)g_i(\hat\beta_1)'$, a
first-order Taylor expansion of $\hat\beta_2$ around the truth $\beta_0$,
\emph{additionally} expanding $\hat\beta_1$ around $\beta_0$, gives
\citep[eq.\ 2.3--2.4]{Windmeijer2005}
\[
  \hat\beta_2-\beta_0 = -A_{\beta_0,W_N(\beta_0)}^{-1}b_{\beta_0,W_N(\beta_0)}
    + D_{\beta_0,W_N(\beta_0)}\left(\hat\beta_1-\beta_0\right) + \bigO_p(N^{-3/2}),
\]
In general $D$'s $j$-th column has three terms \citep[eq.\ 2.3]{Windmeijer2005},
one of which involves $G(\beta_0)=\partial^2\bar g(\beta)/\partial\beta\partial\beta'|_{\beta_0}$
and vanishes whenever $C(\beta)=\partial\bar g(\beta)/\partial\beta'$ is
constant in $\beta$ -- exactly \pkg{}'s case, since its moment conditions are
\emph{linear} ($g_i(\beta)=Z_i'(Y_i-X_i\beta)$, so
$C(\beta)=-\frac1N\sum_iZ_i'X_i$ regardless of $\beta$). Writing, for
brevity,
\[
  K_j(\beta_0) \equiv A^{-1}C(\beta_0)'W_N^{-1}(\beta_0)\left.\frac{\partial W_N(\beta)}{\partial\beta_j}\right|_{\beta_0},
\]
the two remaining terms give
\[
  D_{\cdot j} = -K_j(\beta_0)\,W_N^{-1}(\beta_0)\,C(\beta_0)A^{-1}b_{\beta_0,W_N(\beta_0)}
    \ +\ K_j(\beta_0)\,W_N^{-1}(\beta_0)\,\bar g(\beta_0).
\]
This linear-moment-conditions case is precisely the one
\citet{Windmeijer2005} shows the correction \emph{unambiguously} improves
finite-sample accuracy (his eq.\ 3.1--3.2, specialized to a panel AR(1)-type
model). The corrected variance is
\begin{equation}
  \widehat{\mathrm{var}}_c(\hat\beta_2) = V_2 + D\,\widehat{\mathrm{var}}(\hat\beta_1)\,D'
    + D\,V_2 + V_2\,D'. \label{eq:windmeijer}
\end{equation}
\texttt{GMMEngine.estimate\_two\_step\_robust} implements exactly
\eqref{eq:windmeijer}: \texttt{V1\_robust\_windmeijer} plays the role of
$\widehat{\mathrm{var}}(\hat\beta_1)$ (the first-step's own sandwich
variance, accumulated per cluster consistently with clustering elsewhere in
the engine); \texttt{D\_sum} accumulates
$\partial W_N(\beta)/\partial\beta_j|_{\hat\beta_1}$'s two-cluster-sum
components ($-Z_i'X_{i,j}\cdot(Z_i'\hat\varepsilon_1)'$ and its transpose,
per the linear-moment-condition simplification above); and
\texttt{V2\_robust} assembles \eqref{eq:windmeijer} directly. This is the
variance actually reported whenever \texttt{twostep=True, robust=True} --
the combination most commonly used in applied dynamic-panel work, and the
only one for which naive two-step inference is known to be seriously
unreliable.

\subsection{Small-sample degrees-of-freedom corrections}
\label{sec:qc}

When \texttt{small=True}, \pkg{} rescales variances by finite-sample "qc"
factors, matching \texttt{xtabond2.mata} lines 562--566 exactly:
\[
  V \leftarrow V\times
  \begin{cases}
    \dfrac{N_{\mathrm{obs}}}{N_{\mathrm{obs}}-k} & \text{one-step, non-robust}\\[1.4ex]
    \dfrac{N_{\mathrm{obs}}-1}{N_{\mathrm{obs}}-k}\times q_{\mathrm{cl}} & \text{one-step-robust, two-step (robust or not)}
  \end{cases}
\]
where $k$ is the regressor count (net of the IFE-GMM incidental-parameter
penalty, Section~\ref{sec:dfpca}), and
\[
  q_{\mathrm{cl}} = \begin{cases}
    1 & \texttt{cluster()} \text{ induces a partition different from the panel id}\\
    \dfrac{N_{\mathrm{clusters}}}{N_{\mathrm{clusters}}-1} & \text{clustering induces the same partition as the panel id}
  \end{cases}
\]
-- \texttt{GMMEngine.\_cluster\_qc\_factor}. The second case covers both the
default (no \texttt{cluster()} given, so clustering falls back to the panel
id) and an explicit \texttt{cluster()} that happens to induce the same
partition as the panel id (e.g.\ \texttt{cluster(country\_id)} on a panel
already indexed by \texttt{country\_id}): both leave the moment structure
identical to the no-\texttt{cluster()} case, so both receive the same
correction factor. Only a \texttt{cluster()} that genuinely induces a
different grouping than the panel id drops the factor. This asymmetry
(Mata's \texttt{rows(clusts)} conditional) is confirmed directly from
\texttt{xtabond2.mata} line 564. $N_{\mathrm{obs}}$ itself
is the levels-equation observation count for System GMM (matching Stata's
own $\mathtt{NObs}$, computed from the levels rows) and the differenced-equation
count for Difference GMM.

\section{Diagnostic Tests}
\label{sec:diag}

\subsection{Sargan and Hansen overidentification tests}

The (non-robust) Sargan statistic tests $H_0$: all instruments are valid,
\[
  \mathrm{Sargan} = \frac{\hat\varepsilon_1'Z\,W_1\,Z'\hat\varepsilon_1}{\hat\sigma^2}
    \ \sim\ \chi^2_{\,n_{\mathrm{instr}}-k},
\]
valid only under the same homoskedasticity/no-serial-correlation auxiliary
assumption used to build $W_1$. The Hansen statistic,
$\hat\varepsilon_2'Z\,W_2\,Z'\hat\varepsilon_2\sim\chi^2_{\,n_{\mathrm{instr}}-k}$,
is consistent under the actual (possibly heteroskedastic, clustered, serially
correlated) moment covariance and should be preferred whenever available.
Both lose power as instruments proliferate \citep{Roodman2009b}: with many
instruments, both statistics can fail to reject even a misspecified model, so
a very high $p$-value (e.g.\ $>0.99$) signals over-instrumentation rather
than validity.

\subsubsection{\texorpdfstring{The error-variance scale $\hat\sigma^2$}{The error-variance scale (sigma-hat squared)}}
\label{sec:sig2}

\texttt{xtabond2.mata:201} defines \texttt{ErrorEq}, "a dummy for [the]
equation whose errors to use for $\mathrm{sig2}$ ... transformed eq unless
$h=1$" -- i.e.\ $\hat\sigma^2$ is computed from the transformed
(differenced/FOD) equation's residuals when $h>1$, but from the
\emph{levels}-equation residuals when $h=1$ (consistent with \texttt{h=1}'s
own assumption that the transformed error itself has no differencing-induced
structure):
\[
  \hat\sigma^2 = \frac{1}{d}\cdot\frac{1}{N_{\text{eq}}}\sum \hat\varepsilon_{1,\text{eq}}^2,
  \qquad
  (\text{eq}, d) = \begin{cases}
    (\text{levels},\ 1) & h=1\\
    (\text{diff/FOD},\ 1) & \text{FOD}\\
    (\text{diff},\ 2) & \text{otherwise.}
  \end{cases}
\]
This is exactly \texttt{get\_diagnostics}'s branch on \texttt{self.h}.

\subsection{Difference-in-Sargan/Hansen tests}

Instrument \emph{subset} exogeneity is tested by refitting with that subset
excluded and comparing overidentification statistics:
\[
  \mathrm{Diff} = \mathrm{Stat}_{\mathrm{full}} - \mathrm{Stat}_{\mathrm{restricted}}
    \ \sim\ \chi^2_{\,n_{\mathrm{instr,full}}-n_{\mathrm{instr,restricted}}}
\]
under $H_0$: the excluded subset is valid. \pkg{} tests two subsets
automatically, mirroring \texttt{xtabond2}: the System GMM "GMM instruments
for levels" block (identified as the columns that are zero throughout the
differenced rows and non-constant on the levels rows -- excluding the
automatically-added constant instrument), and a second block covering all
\texttt{iv()}/\texttt{iv\_vars} instruments jointly -- every
standard-instrument variable, across every \texttt{IVStyle} group, tested
together as a single combined subset rather than group by group.
\texttt{diagnostics.py}'s \texttt{compute\_diff\_sargan\_tests} refits a
fresh \texttt{GMMEngine} on the restricted instrument set with
\texttt{check\_rank=False}
(the full model was already identification-checked; a diagnostic refit on a
column \emph{subset} of an already-validated $Z$ does not need a second
full-$Z$ rank SVD).

\subsection{\texorpdfstring{Arellano-Bond AR($\ell$) tests}{Arellano-Bond AR(l) tests}}
\label{sec:artests}

The AR($\ell$) test targets $H_0$: no $\ell$-th order serial correlation in
the \emph{original} idiosyncratic error $\varepsilon$ (not the transformed
one), using $w_{it}=\hat\varepsilon_{i,t-\ell}$ (the transformed-equation
residual lagged $\ell$ periods) and the raw moment
$m_\ell=\sum_{i,t}w_{it}\hat\varepsilon_{it}$. Because $\hat\varepsilon$
depends on the estimated $\hat\beta$ rather than the truth, a naive variance
for $m_\ell$ is invalid; the correctly propagated variance (matching
\texttt{xtabond2.mata::\_ARTests}, implemented in
\texttt{GMMEngine.get\_diagnostics::compute\_ar}) is
\[
  v_\ell = w'Hw + (X'w)'\left[-2VZXA\cdot(Z'Hw) + V_{\mathrm{rob}}\cdot(X'w)\right],
  \qquad z_\ell = \frac{m_\ell}{\sqrt{v_\ell}}\ \overset{a}{\sim}\ N(0,1),
\]
where $VZXA=V\cdot X'Z\cdot W$ is the GMM "bread" and $V_{\mathrm{rob}}$/$V$
are the (possibly robust) parameter covariance matrices -- the extra terms
propagate the estimation uncertainty in $\hat\beta$ into the test. By
construction, $\Delta\varepsilon_{it}=\varepsilon_{it}-\varepsilon_{i,t-1}$
is first-order correlated \emph{even under the null of no correlation in
$\varepsilon$} ($\mathrm{Cov}(\Delta\varepsilon_{it},\Delta\varepsilon_{i,t-1})=-\mathrm{Var}(\varepsilon_{i,t-1})\ne0$),
so \textbf{rejecting AR(1) is the expected, uninformative outcome}; failing
to reject AR(2) or higher is the actual validity requirement for lag-2+
instruments (a rejected AR(2) means the original $\varepsilon$ was itself
serially correlated, invalidating those instruments). \texttt{artests}
generalizes $\ell$ beyond Stata's hardcoded AR(1)/AR(2) to any
$\ell=1,\dots,\ell_{\max}$ (set via \texttt{artests}). \texttt{arlevels} redirects the test to the
\emph{levels}-equation residuals (System GMM only), forcing $H$'s $w'Hw$ term
to the identity ($h=1$-style) regardless of the model's own \texttt{h} --
matching \texttt{xtabond2.mata}'s
\texttt{H = \_H(arlevels?1:h, ...)} -- while the $Z'Hw$ cross-term is
\emph{not} forced (Mata's \texttt{psit = \_H(h,...)} always uses the actual
\texttt{h}).

\subsection{\texorpdfstring{Wald/$F$ test}{Wald/F test}}

Joint significance of all regressors uses
$\mathrm{Wald}=\hat\beta'V^{-1}\hat\beta$ (constant excluded, via a
pseudo-inverse guarding against near-singular $V$), reported as
$\mathrm{Wald}/k_{\mathrm{indep}}\sim F_{k_{\mathrm{indep}},\,\mathrm{df}_{\mathrm{resid}}}$
under \texttt{small=True} or $\mathrm{Wald}\sim\chi^2_{k_{\mathrm{indep}}}$
otherwise.

\section{IFE-GMM: An Original Extension for Interactive Fixed Effects}
\label{sec:ife}

\subsection{Motivation}

Sections~\ref{sec:gmm}--\ref{sec:diag} assume $\varepsilon_{it}$ is
idiosyncratic conditional on $\eta_i$. If the true error instead has an
\emph{interactive fixed effects} (multifactor) structure,
\begin{equation}
  u_{it} = \lambda_i'F_t + \varepsilon_{it}, \label{eq:ife}
\end{equation}
with $\lambda_i$ a unit-specific loading on a common, possibly
time-correlated shock $F_t$ \citep{Bai2009}, and if $F_t$ is correlated with
the regressors (a common macro shock driving both $x_{it}$ and $y_{it}$,
for instance), then $\lambda_i'F_t$ acts as an omitted, regressor-correlated
term that first-differencing/FOD does \emph{not} remove (differencing
eliminates the time-\emph{invariant} $\eta_i$, not the time-\emph{varying}
$F_t$). Standard GMM applied to \eqref{eq:model} in the presence of
\eqref{eq:ife} is therefore inconsistent. Note that \eqref{eq:ife}
\emph{nests} the additive fixed-effects case as the special case $r=1$,
$\lambda_i\equiv1$: it is a strict generalization, not a different model
class \citep[Section 1]{Bai2009}.

\subsection{\texorpdfstring{Identification: the $r^2$ restrictions}{Identification: the r-squared restrictions}}
\label{sec:identification}

$\Lambda F'=\Lambda AA^{-1}F'$ for any invertible $r\times r$ $A$, so
$(\Lambda,F)$ is unidentified without normalizing restrictions. Since an
invertible $r\times r$ matrix has $r^2$ free elements, exactly $r^2$
restrictions are needed \citep[Section 3.1]{Bai2009}. Two standard sets
jointly deliver this and uniquely determine $(\Lambda,F)$ given the product
$\Lambda F'$:
\[
  F'F/T = I_r \quad (r(r+1)/2\text{ restrictions, symmetry}),\qquad
  \Lambda'\Lambda = \text{diagonal} \quad (r(r-1)/2\text{ restrictions}),
\]
summing to $r(r+1)/2+r(r-1)/2=r^2$ exactly. \pkg{}'s IFE-GMM loop
(\texttt{ife.py}) enforces $F'F/T=I_r$ directly by construction (factors are
extracted as the top-$r$ right singular vectors of the residual matrix,
rescaled by $\sqrt{T}$: \texttt{F = Vt[:r,:].T * sqrt(n\_times)}), which is
sufficient for its purposes -- the estimator's target is $\beta$, for which
only the \emph{column space} of $F$ (not the individual loadings) matters,
so the $\Lambda'\Lambda$-diagonal restriction is not separately imposed.

\subsection{\texorpdfstring{Defactoring: the projection matrix $M_F$}{Defactoring: the projection matrix M\_F}}

Given $\beta$, the residual $W_i=Y_i-X_i\beta$ follows the pure factor model
$W_i=F\lambda_i+\varepsilon_i$ \citep[eq.\ 9]{Bai2009}. Concentrating out
$\Lambda=W'F/T$ from the least-squares objective yields, for a candidate
$F$, the projection
\begin{equation}
  M_F = I_T - F(F'F)^{-1}F' = I_T - FF'/T, \label{eq:mf}
\end{equation}
\citep[Section 3.2]{Bai2009} with $M_FW_i=M_F\varepsilon_i$ (the factor
component is exactly annihilated). $F$ itself is recovered as the top-$r$
eigenvectors (scaled by $\sqrt T$) of $\sum_iW_iW_i' = WW'$, i.e.\ via SVD of
the residual matrix -- exactly \texttt{ife.py}'s
\texttt{U, S, Vt = svd(E\_mat); F = Vt[:r,:].T * sqrt(n\_times)}, and
\eqref{eq:mf} is exactly \texttt{M\_F = np.eye(n\_times) - (F @ F.T) / n\_times}.
Bai's own baseline iteration alternates
$\hat\beta(F,\Lambda)=\left(\sum_iX_i'X_i\right)^{-1}\sum_iX_i'(Y_i-F\lambda_i)$
(ordinary least squares, given the factor structure) with re-extracting
$(F,\Lambda)$ from the implied pure-factor residuals.

\subsection{From least squares to profile GMM}

Bai's baseline update for $\beta$ is OLS, valid when $X_{it}$ is (conditional
on the factors) exogenous. In a \emph{dynamic} panel, $X_{it}$ includes
$y_{i,t-1}$, mechanically endogenous even after defactoring -- OLS on
defactored data would reproduce the Nickell-type bias
(Section~\ref{sec:gmm}) all over again. \citet{HongSuJiang2022} generalize
Bai's scheme by replacing the OLS update with a GMM update using
\emph{defactored instruments}: given $F$ (hence $M_F$), project $Y$, $X$,
\emph{and} $Z$ onto $M_F$'s column space and re-run standard GMM
(Section~\ref{sec:onetwostep}) on the projected data. They note (their eq.\
2.8) that substituting the \emph{estimated} loadings $\hat\lambda_i$
directly into a per-unit defactoring is a tempting "naive" alternative, but
complicates the asymptotic analysis; their recommended construction (their
eq.\ 2.9) is instead to defactor with the \emph{projection} $M_F$ applied
uniformly, exactly as in \eqref{eq:mf} -- precisely \texttt{ife.py}'s choice:
\texttt{Y\_proj = Y\_mat @ M\_F}, and likewise for every column of $X$ and
every GMM/IV-instrumented variable, \emph{before} the GMM engine ever sees
the data for that iteration.

\subsection{The iterative algorithm}
\label{sec:ifealgo}

\begin{enumerate}[nosep]
  \item \textbf{Initialize.} $\hat\beta^{(0)}\leftarrow$ plain GMM on the
    raw (non-defactored) data (Section~\ref{sec:onetwostep}).
  \item \textbf{Residuals.} $\hat u_{it}^{(m)} = y_{it}-x_{it}'\hat\beta^{(m)}$,
    reshaped into an $N\times T$ matrix $\hat U^{(m)}$.
  \item \textbf{Factors.} $F^{(m)}\leftarrow\sqrt{T}\times$(top-$r$ right
    singular vectors of $\hat U^{(m)}$); $M_F^{(m)}$ via \eqref{eq:mf}.
  \item \textbf{Project.} $\tilde Y^{(m)}=YM_F^{(m)}$, $\tilde
    X^{(m)}_{\cdot,j}=X_{\cdot,j}M_F^{(m)}$ for every regressor $j$, and
    likewise for every GMM-style/IV-style variable feeding the instrument
    matrix.
  \item \textbf{Re-estimate.} $\hat\beta^{(m+1)}\leftarrow$ GMM on
    $(\tilde Y^{(m)},\tilde X^{(m)},\tilde Z^{(m)})$.
  \item \textbf{Converge?} If $\sum_j|\hat\beta^{(m+1)}_j-\hat\beta^{(m)}_j|<\tau$
    (\texttt{ife\_tol}, default $10^{-5}$), stop; else repeat from step 2, up to
    \texttt{ife\_max\_iter} (default 30) iterations.
\end{enumerate}
Diagnostics (Section~\ref{sec:diag}) are computed only once, on the final
converged fit -- they are unnecessary at intermediate iterations, where only
$\hat\beta^{(m)}$ feeds the next step (an optimization exploited in the
package's performance work, orthogonal to the econometrics documented here).

\paragraph{Missing panel cells.} $\hat U^{(m)}$ can contain undefined
entries for two distinct reasons that leave the same kind of hole: a
genuinely unbalanced panel (an $(i,t)$ observation absent from the data),
or simply a regressor's lag-truncated first observation per unit (present
in the data, but with an undefined residual). Rather than a separate
imputation pass, each such cell is filled from the \emph{previous}
iteration's own rank-$r$ reconstruction of $\hat U$ -- $U_{\cdot,1:r}\,
\mathrm{diag}(S_{1:r})\,V_{\cdot,1:r}'$ from that iteration's SVD, or the
cross-sectional mean on the very first iteration, before any reconstruction
exists -- before the SVD in step 3 above. The outer loop already
alternates between a factor-structure estimate and a data reconstruction
the same way an EM algorithm would, so no nested imputation loop is
needed. \pkg{} prints a one-time, informational console notice the first
time this triggers per fit; it does not alter the structure of the
returned result (Section~\ref{sec:limitations} discusses the caveat this
raises under heavy missingness).

\subsection{Convergence}
\label{sec:convergence}

\citet{HongSuJiang2022} prove a \emph{contraction mapping} property for
exactly the update in Section~\ref{sec:ifealgo}: iterating converges quickly
to the true $\beta$ \emph{provided the starting value $\hat\beta^{(0)}$ is
already consistent under interactive fixed effects}. They obtain such a
starting value via a dedicated nuclear-norm-regularized (NNR) estimator.
\pkg{}'s loop instead starts from plain, non-defactored GMM
(Section~\ref{sec:ifealgo}, step 1) -- \emph{not} guaranteed consistent when
factors are strong, since plain GMM is exactly the estimator IFE-GMM exists
to correct. Citing the iterative estimator they compare against
\citep[Jiang et al.\ 2021, as discussed in][]{HongSuJiang2022}: "the
asymptotic properties of \dots iterative estimators depend on the initial
estimator, the population structure and the iterative steps, and they may
not be consistent if the initial estimators are inconsistent." \pkg{} does
not implement an NNR (or other factor-consistent) starting point; empirical
convergence should be read as a favourable sign, not a proof of correctness,
particularly when factor loadings are large relative to $\varepsilon$'s
scale.

\subsection{Factor-number selection}
\label{sec:factorselect}

When \texttt{r='auto'} is used, \pkg{}'s factor-count routine
(\texttt{numfac.estimate\_num\_factors}, a direct port of Stata's
\texttt{xtnumfac.ado}) computes, from the plain-GMM
residual matrix $E$ ($N\times T$, or its EM-PCA imputation if unbalanced) and
its singular values $S_1\ge\dots\ge S_{\min(N,T)}$, $\mu_j=S_j^2/(NT)$:

\paragraph{Bai \& Ng (2002) $IC_{p2}$.} With
$V(k)=\frac1{NT}\sum_{j>k}\mu_j$ (residual variance after removing $k$
factors) and $V(0)=\overline{E^2}$,
\[
  IC_{p2}(k) = \ln V(k) + k\cdot\frac{N+T}{NT}\ln\big(\min(N,T)\big),
  \qquad \hat r_{IC} = \arg\min_{k\ge1} IC_{p2}(k).
\]

\paragraph{Ahn \& Horenstein (2013) eigenvalue ratio (ER).}
\[
  ER(k) = \frac{\mu_{k-1}}{\mu_k}\quad(k\ge1),
  \qquad
  ER(0) = \frac{V(0)/\ln(\min(N,T))}{\mu_1}\quad(\text{the "\texttt{mockEV}" device}),
\]
\[
  \hat r_{ER}=\arg\max_{k\ge0}\ ER(k).
\]
The \texttt{mockEV} term stands in for a "zeroth eigenvalue", the specific
device that lets $\hat r_{ER}=0$ be a reachable outcome -- concluding "no
factor structure" -- which $IC_{p2}$ cannot do on its own. \pkg{} reports
$\hat r_{ER}$ as the default automatic choice (\texttt{fac\_results['best\_er']}),
alongside $IC_{p2}$ for comparison, matching \texttt{xtnumfac}'s own
reported preference.

\subsection{Split-panel jackknife bias correction}
\label{sec:jackknife}

Any estimator with an $\bigO(1/T)$ asymptotic bias term, $\hat\beta(T) =
\beta_0 + B/T + o_p(1/T)$, can be bias-corrected by exploiting how that bias
scales with $T$: splitting the sample into two halves of length $T/2$ gives
$\hat\beta(T/2)=\beta_0+2B/T+o_p(1/T)$ for each half, so
\[
  2\hat\beta(T) - \tfrac12\left(\hat\beta_{t\le\bar t}(T/2)+\hat\beta_{t>\bar t}(T/2)\right)
    = \beta_0 + \left(2\cdot\tfrac{B}{T}-\tfrac12\cdot2\cdot\tfrac{2B}{T}\right) + o_p(1/T)
    = \beta_0 + o_p(1/T),
\]
cancelling the leading bias term exactly -- the classical
Quenouille/Tukey jackknife idea, applied along the \emph{time} dimension
(appropriate here since the bias is specifically $\bigO(1/T)$, an
incidental-parameter-style artifact of estimating $T$-growing factor
structure alongside $\beta$, not an $\bigO(1/N)$ phenomenon). \citet{DhaeneJochmans2015}
formalize exactly this split-panel jackknife for fixed-effect models with an
incidental-parameter bias, with the $g=2$ case of their general
$g$-way-split estimator (their eq.\ 2.3) being precisely
\begin{equation}
  \hat\beta_{\mathrm{BC}} = 2\hat\beta_{\mathrm{full}}
    - \tfrac12\left(\hat\beta_{t\le\bar t}+\hat\beta_{t>\bar t}\right), \label{eq:jackknife}
\end{equation}
matching \texttt{ife.py}'s \texttt{bias\_correction=True} computation
verbatim, with $\bar t$ the midpoint time period, and each half-panel
estimate obtained by re-running the \emph{entire} IFE-GMM loop
(Section~\ref{sec:ifealgo}) on that half, warm-started from the full-panel
$\hat\beta$ (\texttt{beta\_warm\_start}) to avoid re-solving from a cold
start on a shorter, noisier sub-panel.

\subsection{Asymptotic theory: variance vs.\ bias}
\label{sec:asymptotics}

\citet[Theorem 3.4]{HongSuJiang2022} characterize both the asymptotic
variance and bias of exactly the iterative estimator of
Section~\ref{sec:ifealgo}. Two results matter directly for \pkg{}'s design:

\begin{enumerate}[nosep]
  \item The estimator's asymptotic \emph{variance} coincides, to leading
    order, with plain GMM's variance evaluated at the (unknown) true
    factors -- no additional term for the sampling uncertainty of the
    \emph{estimated} factors is needed. This validates \pkg{}'s practice of
    reporting the plain GMM sandwich/Windmeijer-corrected variance
    (Section~\ref{sec:windmeijer}) computed on the \emph{converged,
    projected} data, with no IFE-specific variance adjustment.
  \item The estimator's asymptotic \emph{bias}, in contrast, has three
    non-zero terms in general (analogous to Bai's own $B$/$C$ bias terms for
    the additive-effects incidental-parameter problem), vanishing only under
    special cases: strictly exogenous instruments, no cross-sectional
    heteroskedasticity, and no serial correlation/heteroskedasticity in
    $\varepsilon_{it}$. These do not hold in general dynamic-panel
    applications (the entire point of instrumenting is to accommodate
    endogenous/predetermined regressors).
\end{enumerate}
Together, these results relocate the practical risk of using IFE-GMM without
bias correction (i.e.\ leaving \texttt{bias\_correction=True} unset): it is a risk to the
\emph{point estimate}, not (to leading order) to the reported standard
error -- exactly the split-panel jackknife of Section~\ref{sec:jackknife}'s
target. \texttt{bias\_correction=True} is therefore recommended whenever
point estimates (not just significance tests) are of interest, at roughly
double the computational cost (two additional half-panel IFE-GMM loops).

\subsection{Incidental-parameter degrees-of-freedom penalty}
\label{sec:dfpca}

Estimating $r$ factors and $N+T$ loadings/factors alongside $\beta$ consumes
degrees of freedom beyond $k$ (the regressor count). \pkg{} applies the
penalty \texttt{df\_pca} $=r(N_{\mathrm{groups}}+T_{\mathrm{span}}-r)$
wherever Stata's own small-sample corrections subtract $k$
(Sections~\ref{sec:qc},~\ref{sec:artests}), matching the standard
incidental-parameter count for \eqref{eq:ife} under the
Section~\ref{sec:identification} restrictions: $N\times r+T\times r$
loading/factor elements, less the $r^2$ identifying restrictions -- i.e.\
exactly the parameter count implied by \citet[p.\ 1235]{Bai2009}'s own
identification argument, re-derived in Section~\ref{sec:identification}.

\section{Validation Methodology}
\label{sec:val}

Two independent evidence sources underlie every claim in this note:

\begin{enumerate}[nosep]
  \item \textbf{Line-by-line audit against \texttt{xtabond2.mata}} (and,
    for IFE-GMM's factor-selection step, \texttt{xtnumfac.ado}) -- the
    actual Mata source of Stata's \texttt{xtabond2}/\texttt{xtnumfac}
    commands (the \texttt{.ado} files are thin dispatch wrappers; all
    matrix algebra lives in \texttt{.mata}), read directly rather than
    inferred from documentation.
  \item \textbf{A real Stata run} of a consolidated do-file, with
    \texttt{xtabond2} actually installed, comparing its log line-by-line,
    coefficient-by-coefficient and test-by-test, against a frozen Python
    reference output (both live in \texttt{stata\_validation/} in the
    project's development repository -- not necessarily included in every
    source distribution).
\end{enumerate}
A third, complementary safety net -- \texttt{tests/test\_reference\_values.py}'s
frozen regression fixtures (development repository) -- does not
independently establish Stata-fidelity (it has no external ground truth)
but does guarantee that any change to the estimator does
not silently alter any \emph{other} already-validated computation; every
change was followed by a full \texttt{pytest} pass before being
considered complete.

\subsection{Unbalanced panels: validation status}
\label{sec:unbalanced}

Both the standard GMM engine and IFE-GMM support panels with missing
$(i,t)$ rows -- interior gaps, late-starting units, and early attrition --
mirroring \texttt{xtabond2.mata}'s own dense-grid (\texttt{Fill}) handling
of unbalanced data, rather than requiring the caller to pre-densify the
panel. Two scenarios were cross-validated against a real Stata run,
matching to 6 decimal places on every coefficient and standard error: an
interior gap, and a late-starting individual. The latter was chosen
specifically because it is the one scenario able to distinguish between
two readings of the level equation's sample-selection rule -- relative to
the panel's \emph{global} $t_{\min}$, or to each unit's \emph{own} first
observed period -- that are indistinguishable on every other validated
specification (every unit there starts at the panel's global $t_{\min}$).
The match confirms the already-implemented "global" reading is the one
Stata itself uses; no change was needed. End-of-panel attrition was
validated internally (finite, well-posed inference, across all three
AR-test weighting branches of Section~\ref{sec:artests}) but has not yet
been independently re-run against Stata for that specific boundary case
(Section~\ref{sec:limitations}).

\section{Known Limitations}
\label{sec:limitations}

\begin{description}[nosep]
  \item[Weights (\texttt{pweight}/\texttt{aweight}/\texttt{fweight})]
    Not implemented. \texttt{xtabond2.mata} shows weights entering at least
    four distinct, non-trivial places -- the $W_1$ accumulation, the
    error-variance scale (via a \texttt{pwe = e1:/sqrt(wt)} device, not a
    simple pre-multiplication), the two-step "meat", and the AR tests -- and
    only the first was verified with the same rigor as the rest of this
    note before the decision was made not to ship a partially-verified
    feature.
  \item[\texttt{cluster()}: single variable only] Stata's multi-way
    \texttt{cluster(var1 var2)} (Cameron-Gelbach-Miller alternating-sum
    clustering) is out of scope; most macro panels cluster on the panel id
    itself in any case.
  \item[\texttt{cluster()} + \texttt{twostep}/\texttt{robust}, instruments
    $>$ clusters] $S_2$ (eq.~\ref{eq:w2}) becomes exactly rank-deficient
    (rank $\le N_{\mathrm{clusters}}$). NumPy's Moore-Penrose pseudo-inverse
    and Mata's \texttt{invsym}-based generalized inverse are both
    \emph{valid} generalized inverses of a singular matrix, but not
    numerically the \emph{same} one -- point estimates (not just SE) can
    diverge from Stata in this specific regime (confirmed: away from it,
    e.g.\ \texttt{cluster()} on the panel id itself, results match exactly).
    Mitigation: keep instrument count $\le$ cluster count (e.g.\ via
    \texttt{collapse=True}) -- the same mitigation Stata's own on-screen
    warning for this situation implies.
  \item[\texttt{GMMStyle}/\texttt{IVStyle}] Stata's further
    \texttt{passthru}, \texttt{mz}, and \texttt{split} suboptions of
    \texttt{gmm()}/\texttt{iv()} are not implemented.
  \item[IFE-GMM convergence] Not guaranteed from an inconsistent starting
    point (Section~\ref{sec:convergence}).
  \item[IFE-GMM point-estimate bias] Present without
    \texttt{bias\_correction=True} (Section~\ref{sec:asymptotics}).
  \item[Unbalanced panels: attrition] Interior gaps and late-starting units
    are cross-checked against a real Stata run to 6 decimal places
    (Section~\ref{sec:unbalanced}); end-of-panel attrition is validated
    internally only (no exception; finite inference across all three
    AR-test weighting branches) but not yet against an independent Stata
    run for that specific case.
  \item[IFE-GMM imputation under heavy missingness] Missing panel cells
    feeding the defactoring loop (Section~\ref{sec:ifealgo}) are filled
    from the model's own evolving rank-$r$ reconstruction -- reliable for
    light-to-moderate missingness, but not independently stress-tested for
    panels with a large missing fraction, where this circularity could
    compound the convergence caveat of Section~\ref{sec:convergence}.
\end{description}

\section{Conclusion}

\pkg{} reproduces Stata's \texttt{xtabond2} matrix algebra and diagnostics
to the precision demonstrated in Section~\ref{sec:val}, with every
remaining discrepancy from Stata specifically isolated to two documented,
narrow regimes (Section~\ref{sec:limitations}) rather than left as an
unexamined risk. Its original IFE-GMM extension is grounded, term for term,
in \citet{Bai2009}'s interactive-fixed-effects model,
\citet{HongSuJiang2022}'s profile-GMM generalization to endogenous
regressors, and \citet{DhaeneJochmans2015}'s bias correction, with the two
places its practical implementation departs from the cleanest available
theory (the starting point for iteration; the default of
\texttt{bias\_correction=False}) documented as open questions for the user
to weigh, not settled by fiat. Together with the companion User Manual, this
note is intended to let a user or reviewer verify any single reported number
back to its exact econometric source.

\bibliographystyle{plainnat}
\bibliography{references}

\end{document}
